\chapter*{Abstract}
\addcontentsline{toc}{chapter}{Abstract}
\markboth{Abstract}{Abstract}
% DRAFT 2026-11-28 (Tayyab's agent). TODO(team): read; fill the R-136 numbers.

Indoor 3D scanning today needs LiDAR hardware or expert photogrammetry. \textbf{BuildScan~3D} turns a video recorded
with an ordinary Android phone (no LiDAR) into a metric 3D Gaussian Splatting (3DGS) model of the room, with room
layout, semantic labels, a mesh and a floor plan, delivered to a web browser.

The phone app records the video with ARCore poses, IMU data and live image-quality guidance. Camera poses come from
Structure-from-Motion with lens self-calibration, scaled and levelled with ARCore; the model is trained with
depth, normal and planar priors, a motion-blur model and flat (PGSR-lite) Gaussians. For walls with no texture, where
SfM fails, a sensor-assisted chain recovers the camera path from the gyroscope, the accelerometer and point tracks
through the whole video. Since November one command makes every choice from the recording itself and goes from the
video to the web viewer output without hand-set options.

On the main test room the model reaches 25.16\,dB PSNR on 79 held-out off-path frames; tape-measured lengths agreed
within 0.19\,\% on average with an earlier model (v3), and the v9 recipe makes the room 1.2 to 1.5\,\% short. On a featureless recording, where normal SfM placed 87 of 5{,}629 frames, the sensor chain
reached 24.25\,dB on held-out frames using only commercially usable networks. The automatic pipeline processed a
textured bedroom end to end in 104 minutes on an RTX~3070; on the featureless recording it chose the sensor chain by
itself but scored 1.4\,dB below the hand-built chain. On the main test room the automatic pipeline gave 25.33\,dB
(aligned) on held-out frames and a room length 1.86\,\% short of the tape, in 106 minutes with no hand-set option.
